% ============================================================================
% Մանկավարժական լուծումներ - Դիսկրետ մաթեմատիկա, Գործնական 8 (Ռեկուրենտ առնչություններ).
% Խնդիրներ 2 և 6: Հայերեն տարբերակ, XeLaTeX (fontspec + Sylfaen).
% Կազմել՝ xelatex-ով երկու անգամ։
% Տերմինաբանությունը՝ glossary_discrete_focused.md / glossary_Apr.csv-ի համաձայն։
% ============================================================================
\documentclass[11pt,a4paper]{article}

% ---- Տառատեսակ (XeLaTeX) ----
\usepackage{fontspec}
\setmainfont{Sylfaen}

% Armenian has no hyphenation patterns under plain XeLaTeX, so long lines cannot
% break and would run into the margin. Let TeX stretch interword spacing as a
% last resort instead of producing overfull lines.
\tolerance=2000
\emergencystretch=3.5em

% ---- Մաթեմատիկա ----
\usepackage{amsmath,amssymb,amsthm}
\usepackage{mathtools}

% ---- Դասավորություն ----
\usepackage[margin=2.3cm]{geometry}
\usepackage{enumitem}
\usepackage{booktabs}
\usepackage{array}

% ---- Գրաֆիկա ----
\usepackage{xcolor}
\usepackage[most]{tcolorbox}
\usepackage{tikz}
\usetikzlibrary{arrows.meta,positioning,calc,fit,backgrounds,shapes.misc}
\usepackage{pgfplots}
\pgfplotsset{compat=1.18}
\usepackage{hyperref}
\hypersetup{colorlinks=true,linkcolor=blue!55!black,urlcolor=blue!55!black}

% ---- Գունապնակ (Հայաստանի դրոշի գույները 3-շարք գրաֆիկների համար) ----
\definecolor{armRed}{HTML}{D90012}
\definecolor{armBlue}{HTML}{0033A0}
\definecolor{armOrange}{HTML}{F2A800}
\definecolor{ink}{HTML}{22303C}

\setlength{\parindent}{0pt}
\setlength{\parskip}{5pt}

% ---- Վանդակներ ----
\newtcolorbox{problembox}[1][]{enhanced, breakable, sharp corners=south,
  colback=armOrange!10, colframe=armOrange!75!black, boxrule=0.8pt, arc=2mm,
  fonttitle=\bfseries\color{white}, coltitle=white,
  attach boxed title to top left={xshift=4mm,yshift=-2.4mm},
  boxed title style={colback=armOrange!85!black, arc=1mm},
  title={Խնդիր #1}}

\newtcolorbox{ideabox}{enhanced, breakable,
  colback=armBlue!6, colframe=armBlue!70!black, boxrule=0.8pt, arc=2mm,
  fonttitle=\bfseries\color{white}, coltitle=white,
  attach boxed title to top left={xshift=4mm,yshift=-2.4mm},
  boxed title style={colback=armBlue!75!black, arc=1mm},
  title={Գաղափարը}}

\newtcolorbox{methodbox}{enhanced, breakable,
  colback=gray!8, colframe=gray!55!black, boxrule=0.7pt, arc=2mm,
  fonttitle=\bfseries, title={$\bigstar$\ Հիշարժան մեթոդ}}

\newtcolorbox{answerbox}{enhanced, sharp corners=south,
  colback=green!8, colframe=green!50!black, boxrule=1pt, arc=2mm, left=4mm,right=4mm}

\newcommand{\answer}[1]{\begin{answerbox}\centering\textbf{Պատասխան՝}\quad #1\end{answerbox}}

% ---- Վերնագիր ----
\title{\vspace{-1.2cm}\textbf{Դիսկրետ մաթեմատիկա - Գործնական 8}\\[2pt]
\large Ռեկուրենտ առնչություններ՝ մանրամասն լուծումներ\\[2pt]
\normalsize\itshape դատողությունները՝ քայլ առ քայլ բացված}
\author{}
\date{}

\begin{document}
\maketitle
\vspace{-1.0cm}

\begin{center}\small\itshape Ընդգրկում է Գործնական 8-ի հանձնարարված խնդիրները՝ 2 և 6։\end{center}

% =====================================================================
\section*{Նախապատրաստում՝ ի՞նչ է ռեկուրենտ առնչությունը և ի՞նչ է պահանջվում}

\textbf{Ռեկուրենտ առնչությունը} սահմանում է հաջորդականություն՝ յուրաքանչյուր անդամ
արտահայտելով \emph{նախորդ} անդամներով, մի քանի \textbf{սկզբնական պայմանների} հետ միասին։ Օրինակ՝
$a_{n}=a_{n-1}+a_{n-2}$-ը $a_0=a_1=1$ պայմաններով Ֆիբոնաչիի կանոնն է։

Այս գործնականում հանդիպում են երկու բոլորովին տարբեր խնդիր, և ստորև բերված երկու խնդիրները հենց
դրանցից մեկական են՝

\begin{itemize}[leftmargin=1.4em,itemsep=2pt]
  \item \textbf{Լուծել} ռեկուրենտ առնչությունը - գտնել $a_n$-ի \emph{փակ բանաձև}, որը չի հղում
        նախորդ անդամներին (Խնդիր 2)։ Գործիքը \emph{բնութագրիչ հավասարումն} է։
  \item \textbf{Կառուցել} ռեկուրենտ առնչություն - \emph{հաշվարկային} հարցը թարգմանել
        ռեկուրենտ առնչության, ապա կիրառել այն (Խնդիր 6)։ Գործիքը խելացի \emph{դեպքերի բաժանումն} է։
\end{itemize}

% =====================================================================
\section*{Խնդիր 2 - Գծային ռեկուրենտ առնչության լուծում}

\begin{problembox}[2]
Լուծել ռեկուրենտ առնչությունը
\[
a_n = 5a_{n-1} - 6a_{n-2},\qquad n>2,
\]
$a_1=2,\ a_2=10$ սկզբնական պայմաններով։ («Լուծել»-ը նշանակում է գտնել $a_n$-ի փակ բանաձևը։)
\end{problembox}

\begin{ideabox}
Նախ՝ անվանենք կառուցվածքը։ Այս ռեկուրենտ առնչությունը \textbf{գծային} է (յուրաքանչյուր $a_i$
հանդես է գալիս առաջին աստիճանով, երբեք ոչ $a_i^2$ կամ $a_ia_j$), \textbf{համասեռ} է (աջ կողմում
ոչինչ չի ավելացվում) և ունի \textbf{հաստատուն գործակիցներ} ($5$-ը և $-6$-ը երբեք չեն փոխվում
$n$-ի հետ)։ Հենց այս տեսակի ռեկուրենտ առնչության համար ճիշտ ենթադրությունը երկրաչափական
հաջորդականությունն է՝ $a_n=r^n$։ Այն գլխից հնարած չէ՝ երեք դիտարկում պարտադրում են այն։

\textbf{(i) Ամենապարզ ռեկուրենտ առնչությունն արդեն երկրաչափական է։} Մեկքայլանի $a_n=b\,a_{n-1}$
կանոնն ամեն քայլում բազմապատկում է նույն $b$ հաստատունով, ուստի
\[
a_n=b\,a_{n-1}=b^2a_{n-2}=\cdots=a_0\,b^{\,n}.
\]
Կրկնվող հաստատուն բազմապատկիչը \emph{հենց} երկրաչափական աճ է։ Մեր երկքայլանի կանոնը պարզապես
երկու այդպիսի աճի խառնուրդ է, ուստի սպասում ենք երկրաչափական բաղադրիչներ։

\textbf{(ii) Երկրաչափականը միակ ձևն է, որտեղ մեկ քայլ ետ գնալը նշանակում է հաստատունով
բազմապատկում։} Եթե $a_n=r^n$, ապա $a_{n-1}=\tfrac1r\,a_n$ և $a_{n-2}=\tfrac1{r^2}\,a_n$։
Հաստատուն գործակիցներով ռեկուրենտ առնչությունը հենց սա է անում՝ այն կապում է անդամը նախորդ
հարևանների հետ ֆիքսված թվերի միջոցով։ Երկրաչափական հաջորդականությունը կառուցված է հենց այդ
քայլին համապատասխանելու համար։

\textbf{(iii) Այն փլուզում է անվերջ շատ հավասարումներ մեկի մեջ։} Ռեկուրենտ առնչությունը պետք է
ճիշտ լինի \emph{ամեն} $n$-ի համար միաժամանակ, ինչը անվերջ շատ պահանջ է։ Տեղադրենք $a_n=r^n$, և
յուրաքանչյուր անդամ ունի ընդհանուր $r^{\,n-2}$ արտադրիչը։ Կրճատենք այն, և $n$ ինդեքսն
ամբողջությամբ անհետանում է՝ թողնելով միակ $r^2=5r-6$ հավասարումը։ Այսպիսով՝ մեկ լավ ընտրված $r$
թիվը բավարարում է կանոնին $3$-րդ, $4$-րդ, $100$-րդ քայլում և այդպես շարունակ՝ միաժամանակ։
\emph{Ահա} ենթադրության աշխատելու իրական պատճառը՝ երկրաչափական ձևը միակն է, որը ստիպում է
հաստատուն գործակիցներով ռեկուրենտ առնչությանն ամեն քայլում ասել ճիշտ նույն բանը։ Համապատասխանող
$r$-երի արժեքները ռեկուրենտ առնչության մեջ թաքնված աճի տեմպերն են։
\end{ideabox}

\textbf{Քայլ 1 - Տեղադրենք ենթադրությունը։}
$a_n=r^n$-ը տեղադրենք $a_n=5a_{n-1}-6a_{n-2}$-ի մեջ՝
\[
r^n = 5r^{\,n-1} - 6r^{\,n-2}.
\]
Բաժանենք ամենափոքր $r^{\,n-2}$ աստիճանի վրա (թույլատրելի է, քանի որ $r=0$-ն տալիս է միայն
տրիվիալ հաջորդականությունը)՝
\[
r^2 = 5r - 6 \quad\Longleftrightarrow\quad \boxed{\,r^2 - 5r + 6 = 0\,}.
\]
Սա \textbf{բնութագրիչ հավասարումն} է։

\textbf{Քայլ 2 - Լուծենք այն։}
\[
r^2-5r+6=(r-2)(r-3)=0 \quad\Longrightarrow\quad r_1=2,\quad r_2=3.
\]
Գտանք \emph{երկու} տարբեր երկրաչափական լուծում՝ $2^n$ և $3^n$։

\textbf{Քայլ 3 - Միավորենք դրանք (վերադրում)։}
Քանի որ ռեկուրենտ առնչությունը գծային է և համասեռ, լուծումների \emph{ցանկացած} կոմբինացիա
նորից լուծում է։ Երկու տարբեր արմատների դեպքում ընդհանուր լուծումն է
\[
a_n = C\cdot 2^n + D\cdot 3^n,
\]
որտեղ $C,D$-ն հաստատուններ են։ Ունենք ճիշտ երկու անհայտ՝ մեր երկու սկզբնական պայմաններին
համապատասխանելու համար, և դա պատահականություն չէ։

\textbf{Ինչու՞ է $C\cdot2^n+D\cdot3^n$-ը \emph{ամեն} լուծում, ոչ թե միայն դրանց մի մասը։}
Ենթադրենք մեկ այլ հաջորդականություն ևս բավարարում է $a_n=5a_{n-1}-6a_{n-2}$-ին և պատահաբար
համընկնում է մեր բանաձևի հետ երկու հաջորդական ինդեքսներում։ Այդ դեպքում նրանք պետք է համընկնեն
\emph{ամենուր}՝ յուրաքանչյուր հաջորդ անդամը որոշվում է նախորդ երկուսով (և կանոնը գործում է նաև
հակառակ ուղղությամբ՝ որոշելով նախորդ անդամները)։ Ուստի լուծումը լիովին որոշվում է ցանկացած
երկու հաջորդական արժեքով, և $C,D$ հաստատունները տալիս են հենց այնքան ազատություն, որքան անհրաժեշտ
է $a_1$-ին և $a_2$-ին համապատասխանելու համար։ Այդ համապատասխանեցումը միշտ հաջողվում է, և
միարժեքորեն, քանի որ երկու երկրաչափական բաղադրիչներն անկախ են՝ ստորև $C,D$-ի համար կազմված
գծային համակարգի որոշիչը $2\cdot 9-3\cdot 4=6\neq 0$ է։ Ընդհանուր դեպքում այս որոշիչը
$r_1 r_2(r_2-r_1)$ է՝ ոչ զրո հենց այն դեպքում, երբ երկու արմատները տարբեր են և ոչ մեկը $0$ չէ՝
այստեղ դրանք $2$ և $3$ են, ուստի այն ապահով կերպով ոչ զրո է։ (Հավասար արմատները կզրոյացնեին
այն՝ այն դեպքը, որը պահանջում է մեթոդի վանդակում նշված ձևափոխված տեսքը։) Երբ $C,D$-ն որոշված
են, բանաձևը \emph{հենց} հաջորդականությունն է։

\textbf{Քայլ 4 - Օգտագործենք սկզբնական պայմանները։}
\[
\begin{cases}
a_1=2: & 2C+3D=2,\\[2pt]
a_2=10: & 4C+9D=10.
\end{cases}
\]
Առաջին հավասարումը բազմապատկենք $2$-ով և հանենք երկրորդից՝
\[
(4C+9D)-(4C+6D)=10-4 \;\Longrightarrow\; 3D=6 \;\Longrightarrow\; D=2,
\]
ապա $2C+3(2)=2\Rightarrow C=-2$։

\textbf{Քայլ 5 - Գրենք փակ բանաձևը։}
\[
a_n=-2\cdot 2^n+2\cdot 3^n = 2\cdot 3^n - 2^{\,n+1}.
\]

\textbf{Քայլ 6 - Միշտ ստուգենք։}
\[
a_1=2\cdot3-2^2=6-4=2\ \checkmark\qquad
a_2=2\cdot9-2^3=18-8=10\ \checkmark
\]
\[
a_3 \stackrel{\text{բանաձև}}{=} 2\cdot27-2^4=38
\quad\text{ընդդեմ}\quad
a_3 \stackrel{\text{ռեկուրենտ}}{=} 5\cdot10-6\cdot2=38\ \checkmark
\]

\answer{$\displaystyle a_n = 2\cdot 3^n - 2^{\,n+1}$}

\medskip
\textbf{Բանաձևը տեսնելը։}
Հաջորդականությունը երկու երկրաչափական «շարժիչների» տարբերությունն է։ $3^n$ շարժիչը (կապույտ)
հաղթում է մրցավազքում՝ $2^{n+1}$ շարժիչը (կարմիր) նվազող ուղղում է, ուստի $a_n$-ը (նարնջագույն)
$n$-ի աճին զուգահեռ ավելի ու ավելի սերտ մոտենում է $2\cdot3^n$-ին։

\begin{center}
\begin{tikzpicture}
\begin{axis}[
  width=12.5cm, height=6.2cm,
  xlabel={$n$}, ylabel={արժեք},
  xtick={1,2,3,4,5}, ymin=0, ymax=520,
  legend cell align=left, legend pos=north west,
  grid=both, grid style={gray!20},
  every axis plot/.append style={thick, mark size=2.2pt},
  tick label style={font=\small}, label style={font=\small}
]
\addplot[armBlue, mark=*] coordinates {(1,6)(2,18)(3,54)(4,162)(5,486)};
\addplot[armRed, mark=square*] coordinates {(1,4)(2,8)(3,16)(4,32)(5,64)};
\addplot[armOrange, very thick, mark=triangle*] coordinates {(1,2)(2,10)(3,38)(4,130)(5,422)};
\legend{$2\cdot3^n$ (գերիշխող), $2^{\,n+1}$ (ուղղում), $a_n=2\cdot3^n-2^{\,n+1}$}
\end{axis}
\end{tikzpicture}
\end{center}

\begin{center}
\small
\begin{tabular}{@{}c|cccccc@{}}
\toprule
$n$ & 1 & 2 & 3 & 4 & 5 & 6\\
\midrule
$a_n$ & 2 & 10 & 38 & 130 & 422 & 1330\\
$a_n/a_{n-1}$ & & $5$ & $3.80$ & $3.42$ & $3.25$ & $3.15$\\
\bottomrule
\end{tabular}
\end{center}

\smallskip
Նայեք ստորին տողին՝ քայլից քայլ աճի տեմպը՝ $a_n/a_{n-1}$, հաստատվում է $3$-ի շուրջ՝ ավելի մեծ
բնութագրիչ արմատը։ Մեծ $n$-երի համար $3^n$ անդամը գերակշռում է $2^{\,n+1}$ անդամին, ուստի
հաջորդականությունը վերջում աճում է գրեթե ճիշտ $3^n$-ի պես։ Սա երկրաչափական ենթադրության հակառակ
ուղղությամբ արդյունք տալն է՝ արմատները, որ գտանք, իրականում այն աճի տեմպերն են, որ ռեկուրենտ
առնչությունն ի սկզբանե թաքցնում էր։

\begin{methodbox}
Երկրորդ կարգի գծային համասեռ $a_n=p\,a_{n-1}+q\,a_{n-2}$ ռեկուրենտ առնչության համար՝
\begin{enumerate}[leftmargin=1.5em,itemsep=1pt]
  \item Գրիր $r^2=pr+q$ բնութագրիչ հավասարումը։
  \item \textbf{Երկու տարբեր արմատ} $r_1\neq r_2$՝ ընդհանուր լուծում $a_n=C r_1^n+D r_2^n$։
  \item \textbf{Կրկնակի արմատ} $r_1=r_2=r$՝ ընդհանուր լուծում $a_n=(C+Dn)\,r^n$
        (լրացուցիչ $n$ արտադրիչ՝ անհրաժեշտ, քանի որ $r^n$-ը մենակ տալիս է միայն մեկ լուծում)։
  \item Գտիր $C,D$-ն երկու սկզբնական արժեքներից, ապա ստուգիր։
\end{enumerate}
\end{methodbox}

% =====================================================================
\section*{Խնդիր 6 - Ռեկուրենտ առնչության կառուցում հաշվելու համար}

\begin{problembox}[6]
$10$ երկարության քանի՞ երկուական տող ($0$-երից և $1$-երից կազմված տողեր) է պարունակում
\textbf{երկու հաջորդական զրո} (այսինքն՝ ինչ-որ տեղ հանդիպում է «$00$» բլոկը)։
\end{problembox}

\begin{ideabox}
«$00$» \emph{պարունակող} տողերն ուղղակիորեն հաշվելը դժվար է՝ բլոկը կարող է տեղավորվել բազմաթիվ
համընկնող տեղերում, և դու ավելորդ կհաշվեիր այն տողերը, որոնք պարունակում են այն մեկից ավելի անգամ։

\textbf{Լրացման միջոցով հաշվում։} Շատ ավելի հեշտ է հաշվել \emph{հակառակը}՝ «$00$»
\textbf{չպարունակող} տողերը, և հանել ընդհանուրից՝
\[
\#(\text{պարունակում է }00) \;=\; \underbrace{2^{10}}_{\text{բոլոր տողերը}} \;-\; \#(\text{խուսափում է }00).
\]
Իսկ «$00$-ից խուսափելը» հենց այն տեսակի հատկություն է, որը լավ կառուցվում է մեկ սիմվոլ առ
սիմվոլ՝ ուստի բավարարում է ռեկուրենտ առնչության։
\end{ideabox}

\textbf{Քայլ 1 - Սահմանենք այն, ինչ իրականում հաշվում ենք։}
Դիցուք
\[
f(n) = \#\{n \text{ երկարության երկուական տողեր՝ առանց երկու կից } 0\text{-ի}\}.
\]

\textbf{Քայլ 2 - Գտնենք ռեկուրենտ առնչություն՝ նայելով \emph{վերջին} սիմվոլին։}
Վերցրու $n$ երկարության ցանկացած «լավ» տող և նայիր, թե ինչով է այն վերջանում։ Կան ճիշտ երկու
փոխբացառող դեպք, և յուրաքանչյուր դեպք կպցված է ավելի կարճ լավ տողի վրա՝

\begin{center}
\begin{tikzpicture}[
  font=\small, >=Stealth,
  blk/.style={draw, rounded corners=2pt, minimum height=8.5mm, align=center},
  bit/.style={draw, rounded corners=2pt, minimum height=8.5mm, minimum width=8mm}
]
% Դեպք 1
\node[anchor=east] at (-0.3,0) {վերջ՝ \textbf{1}};
\node[blk, minimum width=4.4cm, fill=armBlue!10] (g1) at (2.6,0)
  {$n-1$ երկ. լավ տող};
\node[bit, fill=armRed!15, right=1.2mm of g1] (one) {$1$};
\node[anchor=west] at ($(one.east)+(0.5,0)$) {$\rightarrow\ f(n-1)$ տարբերակ};

% Դեպք 2
\node[anchor=east] at (-0.3,-1.5) {վերջ՝ \textbf{0}};
\node[blk, minimum width=3.4cm, fill=armBlue!10] (g2) at (2.1,-1.5)
  {$n-2$ երկ. լավ տող};
\node[bit, fill=armRed!15, right=1.2mm of g2] (one2) {$1$};
\node[bit, fill=armOrange!25, right=1.2mm of one2] (zero2) {$0$};
\node[anchor=west] at ($(zero2.east)+(0.5,0)$) {$\rightarrow\ f(n-2)$ տարբերակ};
\node[font=\scriptsize\itshape, anchor=north, text=red!60!black] at (3.6,-2.05)
  {վերջին $0$-ից առաջ սիմվոլը պետք է լինի $1$, այլապես կստեղծվի «$00$»};
\end{tikzpicture}
\end{center}

\begin{itemize}[leftmargin=1.4em,itemsep=1pt]
  \item \textbf{Վերջանում է $1$-ով}՝ ինչ էլ որ լինի դրանից առաջ՝ կարող է լինել \emph{ցանկացած}
        $n-1$ երկարության լավ տող ($1$ ավելացնելը երբեք չի ստեղծում «$00$»)։ Սա տալիս է
        $f(n-1)$ տող։
  \item \textbf{Վերջանում է $0$-ով}՝ «$00$»-ից խուսափելու համար անմիջապես առջևի սիմվոլը պետք է
        լինի $1$, ուստի տողը վերջանում է «$10$»-ով։ Ինչ էլ որ լինի դրանից առաջ՝ ցանկացած
        $n-2$ երկարության լավ տող է։ Սա տալիս է $f(n-2)$ տող։
\end{itemize}
Երկու դեպքերը երբեք չեն համընկնում և ընդգրկում են ամեն ինչ, ուստի
\[
\boxed{\,f(n) = f(n-1) + f(n-2)\,}.
\]
Սա \textbf{Ֆիբոնաչիի կանոնն} է՝ հաշվարկն ինքն է այն բացահայտում։

\textbf{Քայլ 3 - Բազային դեպքեր (փոքր տողերը հաշվել ձեռքով)։}
\[
\begin{aligned}
n=1:\ \{\,0,\ 1\,\} &\Rightarrow f(1)=2,\\
n=2:\ \{\,01,\ 10,\ 11\,\}\ (\text{ոչ } 00) &\Rightarrow f(2)=3,\\
n=3:\ \{\,010,\ 011,\ 101,\ 110,\ 111\,\} &\Rightarrow f(3)=5.
\end{aligned}
\]
($5=3+2$ նկատեք՝ հաստատելով ռեկուրենտ առնչությունը։)

\textbf{Քայլ 4 - Կիրառենք ռեկուրենտ առնչությունը մինչև $n=10$։}
\begin{center}
\small
\begin{tabular}{@{}c|cccccccccc@{}}
\toprule
$n$    & 1 & 2 & 3 & 4 & 5 & 6 & 7 & 8 & 9 & 10\\
\midrule
$f(n)$ & 2 & 3 & 5 & 8 & 13 & 21 & 34 & 55 & 89 & \textbf{144}\\
\bottomrule
\end{tabular}
\end{center}
Այսպիսով՝ կա $f(10)=144$ հատ $10$ երկարության տող, որ \emph{խուսափում} է «$00$»-ից։

\textbf{Քայլ 5 - Լրացման հաշվարկ։}
\[
\#(\text{պարունակում է } 00) = 2^{10} - f(10) = 1024 - 144 = 880.
\]

\answer{$10$ երկարության $880$ երկուական տող պարունակում է երկու հաջորդական զրո։}

\begin{methodbox}
\textbf{Երկու փոխանցելի հնարք՝}
\begin{enumerate}[leftmargin=1.5em,itemsep=1pt]
  \item \emph{Լրացման միջոցով հաշվում} - երբ «առնվազն մեկ վատ կաղապար» պայմանը խճճված է, հաշվիր
        մաքուր տողերը և հանիր ընդհանուրից։
  \item \emph{Ըստ վերջին սիմվոլի դեպքերի բաժանում} - տեղային արգելված կաղապար ունեցող տողերի
        համար ռեկուրենտ առնչություն կառուցելու համար ճյուղավորիր ըստ տողի վերջի՝ յուրաքանչյուր
        ճյուղ ֆիքսված վերջածանց է կպցնում ավելի կարճ վավեր տողի վրա։
\end{enumerate}
\end{methodbox}

\end{document}
